% =============================================================================
%  Robótica (MEEC, IST) — Lab 3: estimação de trajetória com uma IMU
%  Compilar com:  latexmk -pdf main.tex
%  Figuras em Photos/ (cópias de ../*.png; refazer com os scripts de ../../src)
%  Bibliografia em refs.bib (bibtex, estilo IEEEtran)
% =============================================================================
\documentclass[11pt,a4paper]{article}

\usepackage[utf8]{inputenc}
\usepackage[T1]{fontenc}
\usepackage[portuguese]{babel}
\usepackage[margin=2.5cm]{geometry}
\usepackage{graphicx}
\usepackage{booktabs}
\usepackage{tabularx}
\usepackage{float}
\usepackage{amsmath}
\usepackage{siunitx}
\usepackage{caption}
\usepackage[table]{xcolor}
\usepackage{listings}
\usepackage{enumitem}
\usepackage{titlesec}
\usepackage[most]{tcolorbox}
\usepackage[colorlinks=true, linkcolor=azul, citecolor=verde!85!black,
            urlcolor=azul]{hyperref}

\sisetup{output-decimal-marker={,}, per-mode=symbol}

% ----------------------------------------------------------------------------
% Paleta (a mesma do Lab 1): verde para a estrutura, azul para o detalhe, e
% duas cores de sinal usadas só no comentário final (o que correu bem / mal).
% ----------------------------------------------------------------------------
\definecolor{verde}{HTML}{1B7F3B}
\definecolor{azul}{HTML}{1F3A5F}
\definecolor{bom}{HTML}{1B7F3B}
\definecolor{mau}{HTML}{B03A2E}
\colorlet{tabcab}{verde!16}      % linha de cabeçalho das tabelas
\colorlet{tabgrupo}{azul!8}      % linhas que separam grupos dentro das tabelas

\setlist{nosep, leftmargin=*}
\setlist[itemize]{label=\textcolor{verde}{\textbullet}}

\titleformat{\section}{\Large\bfseries\color{verde}}{\thesection}{0.6em}{}
            [{\color{verde}\titlerule[0.8pt]}]
\titleformat{\subsection}{\large\bfseries\color{azul}}{\thesubsection}{0.6em}{}
\titleformat{\paragraph}[runin]{\normalsize\bfseries\color{azul}}{}{0pt}{}

\captionsetup{labelfont={bf,color=azul}, font=small}

\newtcolorbox{caixaresumo}{colback=azul!4, colframe=azul!45, boxrule=0.5pt,
  arc=2pt, left=8pt, right=8pt, top=6pt, bottom=6pt, parbox=false}
\newtcolorbox{destaque}{colback=verde!7, colframe=verde, boxrule=0pt,
  leftrule=3pt, arc=0pt, left=8pt, right=8pt, top=5pt, bottom=5pt}

\renewcommand{\topfraction}{0.9}
\renewcommand{\textfraction}{0.07}
\renewcommand{\floatpagefraction}{0.8}

\lstset{basicstyle=\ttfamily\footnotesize, breaklines=true,
        frame=single, framesep=4pt, rulecolor=\color{azul!45},
        backgroundcolor=\color{azul!4}, keywordstyle=\color{azul}\bfseries,
        stringstyle=\color{verde!85!black}, commentstyle=\color{black!55},
        columns=flexible, keepspaces=true, showstringspaces=false,
        literate={ç}{{\c c}}1 {ã}{{\~a}}1 {é}{{\'e}}1 {ó}{{\'o}}1 {í}{{\'i}}1}
\renewcommand{\lstlistingname}{List}

% \fig{ficheiro}{largura}: inclui a figura se existir; senão deixa um aviso
\newcommand{\fig}[2]{\IfFileExists{#1}{\includegraphics[width=#2]{#1}}%
  {\fbox{\parbox{#2}{\centering\ttfamily\small figura em falta: \detokenize{#1}}}}}

% \todo{...}: ponto do rascunho por preencher (apagar antes de entregar)
\newcommand{\todo}[1]{\textcolor{mau}{\textbf{[Por fazer: #1]}}}

\newcolumntype{L}{>{\raggedright\arraybackslash}X}
\newcommand{\T}[2]{{}^{#1}T_{#2}}
\newcommand{\vect}[1]{\boldsymbol{#1}}
\newcommand{\med}[1]{\tilde{\vect{#1}}}   % grandeza medida pelo sensor

\title{{\color{verde}\textbf{Trajectory estimation using an IMU}}\\[2mm]
       {\large\color{azul} Robótica --- Laboratório 3}}
\author{{Afonso Cruz (102815)} \and {Eduardo Batista (99931)} \and {Henrique Póvoa (102109)}}
\date{Grupo {nº\textbf{37}} \quad --- \quad 18/12/2026}

\begin{document}
\maketitle

\begin{abstract}
\begin{caixaresumo}
\todo{Resumo, a escrever no fim. Estrutura prevista: (1) o que se fez ---
calibração de uma IMU (acelerómetro e giroscópio) com o robô Lite~6 como
referência; inclinómetro, integração do giroscópio e fusão por
\emph{Covariance Intersection}; estimação de dois gestos. (2) Os três ou
quatro números principais: parâmetros estimados, erro de orientação de cada
método, melhor parâmetro da fusão, deriva da posição prevista e medida.
(3) A conclusão sobre até onde se pode confiar na trajetória estimada.}
\end{caixaresumo}
\end{abstract}

% Cada secção no seu ficheiro, para se poder trabalhar em paralelo sem conflitos.
% =============================================================================
\section{Montagem experimental}
\label{sec:montagem}
% =============================================================================

\paragraph{Sensor e aquisição.} A IMU é uma Grove IMU 9DOF, com um
acelerómetro e um giroscópio de três eixos no mesmo integrado
(ICM-20600~\cite{icm20600}); o magnetómetro não é usado. Está ligada por
I\textsuperscript{2}C a um Arduino Uno, que lê os seis valores e os envia por
porta série para o portátil, uma linha por amostra: o instante (relógio do
Arduino, em microssegundos), a aceleração em m\emph{g} e a velocidade angular
em \si{\degree\per\second}. Guardam-se sempre os dados em bruto; toda a
correção é feita depois.
\todo{taxa de amostragem conseguida e alterações ao \emph{sketch} da cadeira
(velocidade da porta série, gamas dos sensores); se a biblioteca der a
temperatura do integrado, gravá-la também.}

\paragraph{O robô como referência.} Para conhecer a orientação e a posição
verdadeiras do sensor, o conjunto foi montado na pinça do UFactory Lite~6 do
Laboratório~2, com um suporte impresso em 3D que substitui um dos dedos
(Figura~\ref{fig:montagem}). O robô executa movimentos programados e regista a
sua pose a cada \SI{10}{\milli\second}; essa pose é a referência com que se
comparam as estimativas. Os dois registos (robô e IMU) são gravados no mesmo
portátil e alinhados no tempo por um movimento brusco no início de cada
ensaio. \todo{confirmar o método de sincronização e a incerteza que deixa.}

Três limitações desta referência: é a pose que o controlador do robô indica,
e não uma medição externa; os motores acrescentam vibração, pelo que os
ensaios de ruído se fazem com o sensor pousado na mesa; e o robô conhece a
orientação em relação à sua base, não à gravidade, pelo que a inclinação da
mesa é estimada junto com a calibração.

\begin{figure}[H]
\centering
\fig{Photos/montagem.png}{0.7\textwidth}
\caption{\todo{fotografia do sensor montado na pinça, com os eixos do
integrado e os da ferramenta desenhados por cima; ao lado, o suporte.}}
\label{fig:montagem}
\end{figure}

\paragraph{Ensaios.} A Tabela~\ref{tab:ensaios} lista os ensaios e a que
tarefa servem. Todos começam e acabam com o sensor parado durante 5 a
\SI{10}{\second}, e todos foram filmados (ligações no fim do relatório).

\begin{table}[H]
\centering
\caption{Ensaios realizados. \todo{preencher datas, durações e repetições
reais.}}
\label{tab:ensaios}
\small
\begin{tabularx}{\textwidth}{@{}clLc@{}}
\toprule
\rowcolor{tabcab}
 & Ensaio & O que se faz & Tarefas \\
\midrule
A & Parado, longo & sensor na mesa, 10 a 30 min & 2 \\
B & Seis posições & cada eixo para cima e para baixo & 2 \\
B2 & Posições ao acaso & 15 a 20 orientações, paradas & 2 \\
C & Rotação conhecida & uma junta do robô a velocidade constante, nos dois sentidos & 2 \\
T & Temperatura & parado, do frio até estabilizar & 2 \\
D & Inclinações conhecidas & varrimento de inclinação imposto pelo robô & 3, 4, 5 \\
D2 & Inclinar com perturbação & o mesmo, com translação sobreposta & 3, 5 \\
E & Mexer e voltar & regressa à pose inicial & 4 \\
F & Gestos & acenar e soco, pelo robô e à mão & 6 \\
G & Deriva da posição & parado e em trajetória conhecida & 6 \\
\bottomrule
\end{tabularx}
\end{table}
     % Montagem experimental (sensor, aquisição, robô como referência)
% =============================================================================
\section{Tarefa 1 --- Modelos dos sensores}
\label{sec:modelos}
% =============================================================================

Designa-se por $\{s\}$ o referencial do sensor e por $\{w\}$ um referencial
fixo, com $z$ vertical para cima; $R$ é a orientação de $\{s\}$ em $\{w\}$.
As grandezas medidas levam til.

\paragraph{Acelerómetro.} Não mede a aceleração $\vect a$ do sensor, mas a
força específica: a aceleração menos a gravidade
$\vect g = (0;\,0;\,-g)$, vista no referencial do sensor. Parado, mede $1\,g$
para cima. O modelo adotado é
\begin{equation}
\med a = S_a\,R^{T}\,(\vect a - \vect g) + \vect b_a + \vect n_a ,
\label{eq:acc}
\end{equation}
com $\vect b_a$ o desvio (\emph{bias}), $S_a$ uma matriz $3\times3$ com os
fatores de escala na diagonal e o desalinhamento entre eixos fora dela, e
$\vect n_a$ ruído branco de variância $\sigma_a^2$ por eixo.

\paragraph{Giroscópio.} Mede a velocidade angular $\vect\omega$ do sensor,
expressa em $\{s\}$:
\begin{equation}
\med\omega = S_g\,\vect\omega + \vect b_g + \vect n_g ,
\label{eq:gyro}
\end{equation}
com os mesmos três tipos de erro. O desvio $\vect b_g$ não é constante: muda
cada vez que o sensor é ligado e com a temperatura~\cite{titterton2004}.

\paragraph{Montagem.} O integrado não está alinhado com o suporte. A rotação
$R_{ms}$ entre os eixos do sensor e os da ferramenta do robô, e a posição
$\vect r$ do sensor em relação ao eixo de rotação, entram nas medições: um
acelerómetro a uma distância $\vect r$ do eixo mede a mais
$\vect\alpha\times\vect r + \vect\omega\times(\vect\omega\times\vect r)$.

\paragraph{Parâmetros a estimar.} Tabela~\ref{tab:param}.

\begin{table}[H]
\centering
\caption{Parâmetros dos modelos e ensaio de onde se tira cada um.}
\label{tab:param}
\small
\begin{tabularx}{\textwidth}{@{}lLcl@{}}
\toprule
\rowcolor{tabcab}
Parâmetro & Significado & N.º & Ensaio \\
\midrule
$\vect b_a$ & desvio do acelerómetro & 3 & B, B2 \\
$S_a$ & escala (e desalinhamento) do acelerómetro & 3 (+3) & B, B2 \\
$\sigma_a$ & ruído do acelerómetro & 3 & A \\
$\vect b_g$ & desvio do giroscópio & 3 & A, e início de cada ensaio \\
$S_g$ & escala do giroscópio & 3 & C \\
$\sigma_g$ & ruído do giroscópio & 3 & A \\
$R_{ms}$, $\vect r$ & orientação e posição do sensor no suporte & 3 + 3 & B2 e C, no robô \\
\bottomrule
\end{tabularx}
\end{table}

\begin{destaque}
\textbf{Porque é que estes parâmetros importam:} as Tarefas 4 e 6 integram as
medições no tempo. Um desvio $b_g$ no giroscópio dá um erro de orientação
$b_g\,t$; um desvio $b$ na aceleração dá um erro de posição
$\tfrac12\,b\,t^2$. Erros pequenos e constantes deixam de ser pequenos.
\end{destaque}
           % Tarefa 1 - modelos dos sensores
% =============================================================================
\section{Tarefa 2 --- Ensaios e estimação dos parâmetros}
\label{sec:calib}
% =============================================================================

Cada ensaio foi desenhado para que o valor verdadeiro seja conhecido; a
diferença para o valor medido é o erro do sensor.

\subsection{Ruído e desvio do giroscópio (ensaio A)}

Com o sensor parado a velocidade angular é nula, logo a média de cada eixo do
giroscópio é o desvio $\vect b_g$ e a dispersão em torno dela é o ruído. Para
o acelerómetro, o mesmo registo dá $\sigma_a$. A análise de
Allan~\cite{elsheimy2008} separa o ruído branco das variações lentas do
desvio e indica até que duração compensa integrar.
\todo{valores de $\vect b_g$, $\sigma_g$, $\sigma_a$; comparar com a ficha
técnica.}

\paragraph{Filtragem.} \todo{filtro usado (mediana ou passa-baixo), frequência
de corte e porquê; efeito no ruído; cuidado: filtrar atrasa o sinal, por isso
aplica-se o mesmo filtro aos dois sensores ou um filtro sem atraso
(para a frente e para trás), já que tudo é pós-processado.}

\begin{figure}[H]
\centering
\fig{Photos/ruido.png}{0.8\textwidth}
\caption{\todo{ruído dos dois sensores: desvio de Allan, ou sinal em bruto e
filtrado.}}
\label{fig:ruido}
\end{figure}

\subsection{Acelerómetro: desvio e escala (ensaios B e B2)}

Parado, o acelerómetro só sente a gravidade, que é conhecida. Usaram-se dois
métodos.

\paragraph{Seis posições.} Com um eixo alinhado com a vertical, para cima e
depois para baixo, o sensor devia ler $+g$ e $-g$. Das leituras $m^{+}$ e
$m^{-}$:
\begin{equation}
b = \frac{m^{+} + m^{-}}{2}, \qquad s = \frac{m^{+} - m^{-}}{2g}.
\end{equation}

\paragraph{Posições ao acaso.} Em qualquer orientação parada, o módulo da
aceleração corrigida tem de ser $g$. Com $N$ orientações bem distribuídas,
procuram-se os parâmetros que minimizam~\cite{tedaldi2014}
\begin{equation}
\sum_{k=1}^{N}\Bigl(\bigl\lVert S_a^{-1}(\med a_k - \vect b_a)\bigr\rVert - g\Bigr)^2 .
\end{equation}
Não exige alinhar o sensor, mas também não diz como ele está montado: o
módulo é igual para qualquer rotação.

\paragraph{Com o robô.} As mesmas orientações, impostas pelo robô, são
conhecidas. Isso permite verificar o método anterior e estimar a rotação
$R_{ms}$ entre o sensor e a ferramenta, juntamente com a inclinação da mesa.

\begin{figure}[H]
\centering
\fig{Photos/calib_acc.png}{0.75\textwidth}
\caption{\todo{módulo da aceleração em cada posição, antes e depois de
calibrar: os pontos devem passar a cair em $1\,g$.}}
\label{fig:calibacc}
\end{figure}

\subsection{Giroscópio: escala (ensaio C)}

Uma junta do robô roda a velocidade constante conhecida, nos dois sentidos e
a várias velocidades, com o eixo de rotação alinhado sucessivamente com cada
eixo do sensor. A escala é o declive da reta que relaciona a velocidade medida
com a imposta, e o desvio a ordenada na origem.

\begin{figure}[H]
\centering
\fig{Photos/calib_gyro.png}{0.75\textwidth}
\caption{\todo{velocidade angular medida contra a imposta pelo robô, por eixo,
com a reta ajustada.}}
\label{fig:calibgyro}
\end{figure}

\subsection{Efeito da temperatura (ensaio T)}

\todo{desvio do giroscópio e do acelerómetro em função da temperatura do
integrado, do arranque a frio até estabilizar; declive em
\si{\degree\per\second} por \si{\celsius}; comparar com a ficha técnica.}

\subsection{Resultados}

\begin{table}[H]
\centering
\caption{Parâmetros estimados (média e dispersão de 3 repetições).
\todo{preencher.}}
\label{tab:calib}
\small
\begin{tabular}{llccc}
\toprule
\rowcolor{tabcab}
Parâmetro & Método & $x$ & $y$ & $z$ \\
\midrule
$b_a$ [m\emph{g}] & seis posições & & & \\
                  & posições ao acaso & & & \\
                  & robô & & & \\
$s_a$ [--]        & seis posições & & & \\
                  & posições ao acaso & & & \\
                  & robô & & & \\
$b_g$ [\si{\degree\per\second}] & parado & & & \\
$s_g$ [--]        & robô & & & \\
$\sigma_a$ [m\emph{g}], $\sigma_g$ [\si{\degree\per\second}] & parado & & & \\
\bottomrule
\end{tabular}
\end{table}

\todo{comentário aos resultados: os três métodos concordam? qual é a
dispersão entre repetições? quanto mudou o desvio do giroscópio entre sessões?
orientação do sensor no suporte.}
           % Tarefa 2 - ensaios e parâmetros
% =============================================================================
\section{Tarefa 3 --- Inclinómetro com o acelerómetro}
\label{sec:incl}
% =============================================================================

Se o sensor não acelera, o acelerómetro calibrado mede só a gravidade, e a
sua direção no referencial do sensor dá a inclinação. Com
$\vect a = (a_x;\,a_y;\,a_z)$ a medição corrigida,
\begin{equation}
\phi = \operatorname{atan2}(a_y,\ a_z), \qquad
\theta = \operatorname{atan2}\!\bigl(-a_x,\ \sqrt{a_y^2 + a_z^2}\bigr),
\label{eq:incl}
\end{equation}
onde $\phi$ é o rolamento e $\theta$ a arfagem. A rotação em torno da
vertical não altera a gravidade vista pelo sensor, por isso não se obtém.

O método tem dois limites. É ruidoso, porque cada estimativa usa uma só
amostra. E assume aceleração nula: qualquer aceleração do sensor é lida como
inclinação. O módulo da medição indica quando isso acontece, porque se afasta
de $1\,g$.

\todo{resultados no ensaio D (erro médio e máximo contra o ângulo do robô) e
no D2 (erro durante a perturbação); remeter para a Figura~\ref{fig:fusao}.}
           % Tarefa 3 - inclinómetro
% =============================================================================
\section{Tarefa 4 --- Orientação com o giroscópio}
\label{sec:gyro}
% =============================================================================

A orientação obtém-se integrando a velocidade angular corrigida. Para evitar
a singularidade dos ângulos de Euler (a \ang{90} de arfagem dois eixos
alinham-se, como na singularidade do pulso do robô no Laboratório~2), a
orientação é guardada num quaternião unitário $q$ e integrada
por~\cite{titterton2004}
\begin{equation}
\dot q = \tfrac12\, q \otimes (0;\ \vect\omega), \qquad
q_{k+1} = q_k \otimes \exp\!\bigl(\tfrac12\,\vect\omega_k\,\Delta t\bigr),
\label{eq:quat}
\end{equation}
com normalização a cada passo. Os ângulos só são calculados no fim, para
apresentação. A orientação inicial vem do inclinómetro, com o sensor parado.

O método responde depressa e não é afetado pela aceleração do sensor, mas o
erro cresce com o tempo: um desvio residual $\delta b_g$ acumula
$\delta b_g\,t$, e o ruído soma-se como um passeio aleatório.

\todo{resultados no ensaio D (erro ao longo do tempo) e no E (erro ao
regressar à pose inicial; deriva em \si{\degree\per\minute}), com e sem a
calibração da Tarefa~2.}
           % Tarefa 4 - orientação com o giroscópio
% =============================================================================
\section{Tarefa 5 --- Fusão das duas estimativas}
\label{sec:fusao}
% =============================================================================

As duas estimativas de inclinação têm defeitos complementares: a do
acelerómetro não deriva mas é ruidosa e falha quando o sensor acelera; a do
giroscópio é suave e imune à aceleração, mas deriva.

\subsection{Covariance Intersection}

Sejam $\hat{\vect x}_a$ e $\hat{\vect x}_g$ as duas estimativas de
$(\phi;\,\theta)$, com covariâncias $P_a$ e $P_g$. A
\emph{Covariance Intersection}~\cite{julier1997} combina-as com um peso
$w \in [0,1]$:
\begin{equation}
P^{-1} = w\,P_a^{-1} + (1-w)\,P_g^{-1}, \qquad
\hat{\vect x} = P\bigl(w\,P_a^{-1}\hat{\vect x}_a + (1-w)\,P_g^{-1}\hat{\vect x}_g\bigr).
\label{eq:ci}
\end{equation}
Com $w=1$ fica só o acelerómetro e com $w=0$ só o giroscópio. Ao contrário
de um filtro de Kalman, a CI não assume que os erros das duas estimativas são
independentes, e o resultado nunca é mais confiante do que os dados
justificam. Aqui isso é relevante: a estimativa fundida de cada instante é o
ponto de partida da integração seguinte, e os erros do acelerómetro (desvio
residual, acelerações que duram várias amostras) estão correlacionados no
tempo.

\todo{como se obtêm $P_a$ (ruído do acelerómetro propagado
por~\eqref{eq:incl}; aumentada quando o módulo se afasta de $1\,g$) e $P_g$
(cresce a cada passo com $\sigma_g$); se $w$ é fixo ou escolhido a cada
instante para minimizar o traço de $P$.}

\subsection{Resultados}

\begin{figure}[H]
\centering
\fig{Photos/fusao_tempo.png}{\textwidth}
\caption{\todo{um varrimento de inclinação (ensaios D e D2): ângulo imposto
pelo robô, estimativa do acelerómetro, do giroscópio e da fusão; por baixo, o
erro de cada uma.}}
\label{fig:fusao}
\end{figure}

\begin{figure}[H]
\centering
\fig{Photos/fusao_w.png}{0.7\textwidth}
\caption{\todo{erro quadrático médio da fusão em função de $w$, com o mínimo
assinalado; estimativas no tempo para três valores de $w$.}}
\label{fig:ciw}
\end{figure}

\begin{table}[H]
\centering
\caption{Erro de inclinação de cada método contra a referência do robô.
\todo{preencher.}}
\label{tab:fusao}
\small
\begin{tabular}{lcccc}
\toprule
\rowcolor{tabcab}
 & \multicolumn{2}{c}{Ensaio D} & \multicolumn{2}{c}{Ensaio D2 (com perturbação)} \\
\rowcolor{tabcab}
Método & RMS & máx. & RMS & máx. \\
\midrule
Acelerómetro (Tarefa 3) & & & & \\
Giroscópio (Tarefa 4)   & & & & \\
CI, melhor $w$          & & & & \\
Filtro de Kalman        & & & & \\
\bottomrule
\end{tabular}
\end{table}

\subsection{Comparação com um filtro de Kalman}

\todo{opcional, a confirmar com o docente. As Tarefas 3 e 4 são as duas
metades de um filtro de Kalman~\cite{kalman1960}: o giroscópio faz a previsão
e o acelerómetro a correção. Estado: orientação e desvio do giroscópio.
Comparar com a CI nos mesmos dados (Tabela~\ref{tab:fusao}) e discutir a
confiança de cada um (covariância estimada contra erro real).}
           % Tarefa 5 - fusão (CI)
% =============================================================================
\section{Tarefa 6 --- Gestos e estimação da trajetória}
\label{sec:gestos}
% =============================================================================

\subsection{Da orientação à posição}

Com a orientação $R(t)$ da Tarefa~5, a aceleração do sensor em $\{w\}$
obtém-se de~\eqref{eq:acc}, e a posição integrando duas vezes:
\begin{equation}
\vect a = R\,S_a^{-1}(\med a - \vect b_a) + \vect g, \qquad
\vect v(t) = \int_0^t \vect a\,d\tau, \qquad
\vect p(t) = \int_0^t \vect v\,d\tau .
\label{eq:pos}
\end{equation}

\begin{destaque}
\textbf{Previsão.} Um erro constante $b$ na aceleração dá um erro de posição
$\tfrac12\,b\,t^2$. Há duas fontes: o desvio residual do acelerómetro, e o
erro de inclinação $\delta\theta$, que deixa por subtrair $g\sin\delta\theta$.
Com \ang{1} de erro de inclinação são \SI{0.17}{\meter\per\second\squared},
ou \SI{77}{\centi\meter} ao fim de \SI{3}{\second}. Por isso os gestos são
curtos e começam e acabam com o sensor parado.
\end{destaque}

Como cada gesto começa e acaba em repouso, a velocidade nos dois extremos é
conhecida (zero). \todo{correção usada: retirar à velocidade a deriva linear
entre os dois instantes de repouso; e, tudo sendo pós-processado, um
suavizador~\cite{rauch1965} em vez de um filtro só para a frente.}

\subsection{Validação com o robô (ensaios F e G)}

Antes dos gestos à mão, o método foi aplicado a dois movimentos executados
pelo robô, de que se conhece a posição em cada instante: um ``acenar'' e um
``soco'', desenhados nas juntas como o swing do Laboratório~2.

\begin{figure}[H]
\centering
\fig{Photos/deriva.png}{0.75\textwidth}
\caption{\todo{erro de posição ao longo do tempo, com o sensor parado e numa
trajetória do robô, com a curva prevista $\tfrac12 b t^2$; com e sem
calibração e com e sem correção de velocidade nula.}}
\label{fig:deriva}
\end{figure}

\todo{erro de posição no fim de cada gesto do robô; até que duração o erro
fica abaixo de um valor útil (por exemplo \SI{5}{\centi\meter}).}

\subsection{Gestos à mão}

Com o melhor $w$ da Tarefa~5 e o sensor seguro na mão, fizeram-se devagar um
gesto de acenar e um de soco, \todo{$n$} vezes cada.

\begin{figure}[H]
\centering
\fig{Photos/gestos.png}{\textwidth}
\caption{\todo{acenar e soco à mão: orientação ao longo do tempo e trajetória
estimada, com fotogramas do vídeo; referência geométrica (comprimento da
régua, raio do antebraço).}}
\label{fig:gestos}
\end{figure}

\todo{resultados: amplitude e forma de cada gesto contra a referência;
repetibilidade entre repetições; o que se reconhece e o que não.}
           % Tarefa 6 - gestos e trajetória
% =============================================================================
\section{Para além do enunciado: controlar o robô com a mão}
\label{sec:extra}
% =============================================================================

\todo{secção opcional. A fusão da Tarefa~5 a correr em tempo real: o sensor
preso às costas da mão comanda o robô. (i) Primeiro no modelo MuJoCo do
Lite~6 do Laboratório~2; (ii) depois no robô, com poucas juntas, gamas e
velocidades limitadas, filtro e uma tecla de segurança. Dois modos: a pinça
copia a inclinação da mão; e ``inclinar para mover''. Mostrar também que
seguir diretamente a posição integrada não funciona (Secção~\ref{sec:gestos}).
Uma figura e a ligação para o vídeo.}
        % Para além do enunciado: controlar o robô com a mão
% =============================================================================
\section{Comentário aos resultados}
\label{sec:comentario}
% =============================================================================

\subsection*{\textcolor{bom}{Pontos bons}}

\begin{itemize}[itemsep=2pt]
  \item \todo{a preencher com resultados: o que a referência do robô permitiu
        afirmar; concordância entre métodos de calibração; previsão e medição
        da deriva.}
\end{itemize}

\subsection*{\textcolor{mau}{O que não correu tão bem}}

\begin{itemize}[itemsep=2pt]
  \item \todo{limitações reais encontradas: taxa de amostragem, vibração do
        robô, guinada sem referência, referência que é a do controlador.}
\end{itemize}

\subsection*{\textcolor{azul}{O que podíamos melhorar}}

\begin{itemize}[itemsep=2pt]
  \item \todo{por exemplo: magnetómetro para a guinada; medição externa da
        posição; compensação de temperatura; amostragem mais rápida.}
\end{itemize}

% =============================================================================
\section*{Contribuição dos membros}
\todo{a acordar entre os três.}

\section*{Declaração de uso de IA}
\todo{a escrever pelo grupo, de acordo com o que a cadeira pede.}
   % comentário, contribuições, declaração de IA
% =============================================================================
\section*{Vídeos dos ensaios}
% =============================================================================

\todo{ligações. O enunciado pede vídeos dos ensaios da Tarefa~2; juntar os
dos gestos da Tarefa~6 e o do controlo do robô.}
       % vídeos

% =============================================================================
\bibliographystyle{IEEEtran}
\bibliography{refs}

\end{document}
